% Solution: Gibbons (1992) Problem 2.19. Input by \solitem. Check: check_T2.py (item 11).
\begin{exsol}{Gibbons Problem 2.19}{137}{Exam level}{3}
\asked In the three-period Rubinstein model of Section 2.1.D, whose backward-induction outcome we solved, what
is the subgame-perfect Nash equilibrium?

\method Gibbons \S2.1.D (pp.~68--70). The pie is $1$ and the discount factor $\delta$. Period~1: player~1
proposes $(s_1, 1-s_1)$ and player~2 accepts or rejects. Period~2: player~2 proposes $(s_2, 1-s_2)$ and
player~1 accepts or rejects. Period~3: player~1 receives $s$ and player~2 receives $1-s$, with $0 < s < 1$.
The backward-induction outcome is the \emph{path}; the SPNE is the complete strategy of each player, that is, an
action at every history, including those that never occur.

\solution
\textbf{Backward induction.} In period~2 player~1 accepts $s_2$ iff $s_2 \ge \delta s$ (rejecting gives him $s$
one period later). Player~2 proposes $s_2^* = \delta s$ and keeps $1 - \delta s$. In period~1 player~2
accepts $s_1$ iff $1 - s_1 \ge \delta(1 - s_2^*) = \delta(1 - \delta s)$. Player~1 proposes
\[
  s_1^* = 1 - \delta(1 - \delta s) = 1 - \delta + \delta^2 s ,
\]
and it is accepted: payoffs $\bigl(1 - \delta + \delta^2 s,\ \delta - \delta^2 s\bigr)$.

\textbf{The SPNE.}
\begin{itemize}
  \item \emph{Player~1:} in period~1 propose $s_1^* = 1 - \delta(1 - \delta s)$; in period~2, after any
  proposal $s_2$ by player~2, accept iff $s_2 \ge \delta s$.
  \item \emph{Player~2:} in period~1, after any proposal $s_1$, accept iff $1 - s_1 \ge \delta(1 - \delta s)$,
  i.e.\ iff $s_1 \le s_1^*$; in period~2 propose $s_2^* = \delta s$.
\end{itemize}
Both strategies specify behaviour after every proposal, not only the equilibrium one. This is what makes it an
SPNE rather than just a description of the outcome.

\textbf{Verification (one-stage deviations).} \emph{Period~2.} If player~1 rejects $s_2 \ge \delta s$
he gets $\delta s \le s_2$ in value. If player~2 offers $s_2 < \delta s$ it is rejected and she gets
$\delta(1 - s) < 1 - \delta s$ (the last inequality is $1 - \delta s - \delta + \delta s = 1 - \delta > 0$), and
if she offers more than $\delta s$ she gives away surplus. \emph{Period~1.} Player~2 compares accepting
$1 - s_1$ with rejecting and getting $\delta(1 - \delta s)$ in value. Player~1 compares demanding $s_1 \le
s_1^*$ (accepted, payoff $s_1 \le s_1^*$) with demanding more (rejected, payoff $\delta^2 s$). Since
$s_1^* - \delta^2 s = 1 - \delta > 0$, the best demand is $s_1^*$. At all histories the strategies are best
replies, so the profile is subgame perfect.

\interp The backward-induction outcome only says that agreement occurs immediately at $(s_1^*, 1 - s_1^*)$.
The SPNE adds the threats and acceptance rules that make this outcome optimal, and they hold also after
off-path offers. Such rules are the content of Rubinstein's argument (item~19 verifies the stationary
version in the infinite horizon).

\verify \texttt{check\_T2.py} (item~11) tests the formulas for three parameter pairs ($\delta = 0.9$, $s = 0.5$;
$\delta = 0.5$, $s = 0.3$; $\delta = 0.8$, $s = 0.9$): $s_2 = \delta s$, $s_1 = 1 - \delta + \delta^2 s$, player~2 indifferent at $s_1^*$, and on a grid
of step $\frac1{1000}$ no demand by player~1 beats $s_1^*$ given the acceptance rule. Gibbons prints no answer.

\answer{Player~1: propose $s_1^* = 1 - \delta(1 - \delta s)$ in period~1; in period~2 accept iff $s_2 \ge \delta s$.
Player~2: in period~1 accept iff $1 - s_1 \ge \delta(1 - \delta s)$; in period~2 propose $s_2^* = \delta s$. On the
path player~2 accepts $s_1^*$ and the payoffs are $(1 - \delta + \delta^2 s,\ \delta - \delta^2 s)$.}
\end{exsol}
